% ============================================================
% 付録案A：漸化式の発展パターン
% ============================================================
% このファイルは現行本文には接続していない草案である.
% 「特殊な公式を覚える」のではなく,漸化式を見たときに
% どの方向へ発展させればよいか,という思考の型を主題にする.

\chapter{付録案A：漸化式の発展パターン}

\section*{この案の狙い}

漸化式を学ぶと,さまざまな形の問題に出会う.
一次漸化式,二次漸化式,分数式を含む漸化式,差分を使う漸化式.
一見すると,形が変わるたびに別の解法を覚えなければならないように見える.

しかし,少し遠くから見ると,何度も現れる「発展の型」がある.

\begin{itemize}
  \item 高階の漸化式を,状態を増やして一段の規則に変える.
  \item 値そのものではなく,固定点からのずれや逆数を新しい変数にする.
  \item 毎回変わる値の中から,変化しない量を探す.
  \item 正確な一般項ではなく,成長率やスケールを先に調べる.
  \item 一回の更新を行列や関数として記録し,更新の積全体を調べる.
  \item 刻み幅を細かくして,漸化式を微分方程式へ近づける.
\end{itemize}

ここでいう「パターン」は,数列の項に現れる模様ではない.
問題にぶつかったとき,次にどんな見方を試せばよいかという,数学的な発展の型である.

この案では,一つの公式を最後まで解き切ることよりも,
「なぜその変形を思いつくのか」を明らかにする.
漸化式を解くことが目的だったはずなのに,いつの間にか行列,力学系,
不変量,微分方程式,漸近解析へつながっていく.
その広がりそのものを,一つの物語として読むことを目指す.

\section{局所的な規則から,全体の構造へ}

\subsection{漸化式は短いが,問いは長い}

漸化式
\[
  a_{n+1}=f(a_n),\qquad a_0=\alpha
\]
は,次の項を作る方法だけを指定している.
この意味で,漸化式は局所的な規則である.

一方,私たちが知りたいことはしばしば大域的である.

\begin{itemize}
  \item $a_n$ は有限の値に近づくのか.
  \item 近づくなら,その値は何で決まるのか.
  \item 初期値を少し変えると,遠い将来も似たままなのか.
  \item $a_n$ の大きさは $n$ に対してどの程度の速さで増えるのか.
  \item 一般項が書けなくても,周期や発散を判定できるのか.
\end{itemize}

したがって,漸化式を見たら,ただちに一般項を求めようとするだけでは足りない.
まず,何を状態として持ち,どの量が変化し,どの量が保存されるのかを考える.
これが以下のパターンの出発点である.

\subsection{同じ漸化式に,複数の問題が隠れている}

例えば
\[
  a_{n+1}=2a_n+1,\qquad a_0=0
\]
を考える.
一般項は $a_n=2^n-1$ である.
しかし,この一行から少なくとも三つの異なる問題が生まれる.

\begin{enumerate}
  \item 項を正確に求める問題.
  \item $a_n$ がどれくらいの速さで増えるかを調べる問題.
  \item 初期値や係数を変えたとき,同じ構造が残るかを調べる問題.
\end{enumerate}

一つ目には解法がある.
二つ目には成長率という言葉がある.
三つ目には,変数変換,行列,安定性という言葉が現れる.
このように,問題を少し変えると,必要な数学も変わる.

\section{発展パターン1：高階の規則を一段の状態更新にする}

\subsection{二つ前まで必要なら,状態を二次元にする}

二階漸化式
\[
  a_{n+2}=p a_{n+1}+q a_n
\]
では,\,$a_{n+1}$ を決めるために $a_n$ だけではなく,その一つ前の値も必要になる.
一段の規則ではないから扱いにくい,と思うかもしれない.

そこで,値を一つではなく,二つ組にして持つ.
\[
  X_n=
  \begin{pmatrix}
    a_{n+1}\\
    a_n
  \end{pmatrix}.
\]
すると,
\[
  X_{n+1}
  =
  \begin{pmatrix}
    a_{n+2}\\
    a_{n+1}
  \end{pmatrix}
  =
  \begin{pmatrix}
    p a_{n+1}+q a_n\\
    a_{n+1}
  \end{pmatrix}
  =
  \begin{pmatrix}
    p&q\\
    1&0
  \end{pmatrix}
  X_n
\]
となる.
行列
\[
  A=
  \begin{pmatrix}
    p&q\\
    1&0
  \end{pmatrix}
\]
を用いれば,
\[
  X_{n+1}=AX_n,\qquad X_n=A^nX_0
\]
である.

ここで起きたことは,単なる記号の置き換えではない.
二階漸化式を「二次元の空間の中を一段ずつ動く規則」に変えたのである.
高階の漸化式でも,必要な過去の値をすべて状態に入れれば同じことができる.

\subsection{次数ではなく,状態の次元を見る}

$k$ 階の線形漸化式
\[
  a_{n+k}=c_{k-1}a_{n+k-1}+\cdots+c_1a_{n+1}+c_0a_n
\]
に対して,
\[
  X_n=
  \begin{pmatrix}
    a_{n+k-1}\\
    a_{n+k-2}\\
    \vdots\\
    a_n
  \end{pmatrix}
\]
と置くと,
\[
  X_{n+1}=AX_n
\]
という $k$ 次元の行列の漸化式になる.

この見方をすると,漸化式の「階数」は,過去の情報を何個保存すれば未来が決まるかを表す.
二階だから難しいのではない.
未来を決めるために,状態を二次元で持つ必要があるのである.

\subsection{行列の積は,規則の繰り返しを記録する}

係数が一定でない漸化式
\[
  X_{n+1}=A_nX_n
\]
では,
\[
  X_n=A_{n-1}A_{n-2}\cdots A_1A_0X_0
\]
となる.
行列は一般には掛ける順序を交換できないので,
「どの規則をどの順番で適用したか」がそのまま残る.

これは,漸化式を単なる数の計算から,
操作の合成を研究する問題へ変える.
例えば,二つの行列 $A,B$ を交互に使うなら,
\[
  X_{2m}=(BA)^mX_0
\]
となる.
一回ずつの動きは複雑でも,二回分を一つの行列とみなせば,
周期的な規則の長期的な挙動を調べられる.

\section{発展パターン2：固定点を中心に座標を取り直す}

\subsection{なぜ固定点が中心になるのか}

漸化式
\[
  a_{n+1}=f(a_n)
\]
で,ある値 $L$ が
\[
  f(L)=L
\]
を満たすとき,\,$L$ を固定点という.
もしある時点で $a_n=L$ なら,その後もずっと $L$ のままである.

数列が $L$ に近づくかを調べたいなら,値そのものよりも,
\[
  e_n=a_n-L
\]
という誤差を見る方が自然である.
固定点のまわりでは,
\[
  e_{n+1}=f(L+e_n)-L
\]
となる.
ここで $e_n$ が小さいとき,関数の変化を一次近似すれば,
\[
  e_{n+1}\simeq f'(L)e_n
\]
である.
したがって,局所的には $f'(L)$ が誤差を何倍するかを決める.

\begin{itemize}
  \item $|f'(L)|<1$ なら,誤差は縮み,固定点は安定である.
  \item $|f'(L)|>1$ なら,誤差は広がり,固定点は不安定である.
  \item $|f'(L)|=1$ なら,一次近似だけでは結論が出ない.
\end{itemize}

「収束するか」という問いが,「固定点からの誤差が縮むか」という問いに変わった.
これが座標変換の力である.

\subsection{一次漸化式は,固定点を原点に移すと終わる}

\[
  a_{n+1}=ra_n+b
\]
を考える.
固定点 $L$ は
\[
  L=rL+b
\]
を満たすので,\,$r\neq1$ なら
\[
  L=\frac{b}{1-r}
\]
である.
$e_n=a_n-L$ とおくと,
\[
  e_{n+1}
  =a_{n+1}-L
  =ra_n+b-L
  =r(a_n-L)
  =re_n.
\]
つまり,定数項 $b$ は固定点を原点へ移すことで消える.
漸化式を解いたというより,問題の中心を選び直したのである.

\subsection{逆数を取ると,非線形が一次になることがある}

次の漸化式を考える.
\[
  a_{n+1}=\frac{a_n}{1+a_n},
  \qquad a_0>0.
\]
このままでは分数が毎回現れる.
しかし,\,$a_n$ が正なので逆数を考えることができる.
\[
  \frac1{a_{n+1}}
  =\frac{1+a_n}{a_n}
  =\frac1{a_n}+1.
\]
ここで
\[
  b_n=\frac1{a_n}
\]
とおけば,
\[
  b_{n+1}=b_n+1
\]
となる.
したがって,
\[
  b_n=b_0+n,\qquad
  a_n=\frac1{n+\frac1{a_0}}.
\]

分数式の漸化式を解くために,分母を払う計算をするのではなく,
「どの量なら毎回同じ増え方をするか」を探した.
変数変換は,式を飾るための技巧ではなく,規則を単純化するための座標の選択である.

\subsection{分数一次変換は,固定点の比に変えると等比になる}

より一般に,
\[
  a_{n+1}=\frac{\alpha a_n+\beta}{\gamma a_n+\delta}
\]
を考える.
分母が $0$ にならない範囲で,これは分数一次変換と呼ばれる.
固定点 $\lambda,\mu$ があるとき,
\[
  b_n=\frac{a_n-\lambda}{a_n-\mu}
\]
という比を作ると,ある定数 $q$ に対して
\[
  b_{n+1}=q b_n
\]
となる.

値そのものは分数式で動くのに,
二つの固定点からのずれの比は等比数列になる.
これは「複雑な規則の中に,適切な座標を取ると単純な規則が隠れている」
という発展パターンの典型例である.

\section{発展パターン3：変化しない量を探す}

\subsection{解を求める前に,保存量を探す}

漸化式の各項が複雑に変化していても,
項を組み合わせると一定になることがある.
そのような量を不変量,または保存量という.

微分方程式や物理では,エネルギーや運動量が保存されることで運動の全体像が見える.
漸化式でも同じ発想が使える.
一般項を求める代わりに,軌道がどの曲線上を動くかを特定するのである.

\subsection{振動する二階漸化式の不変量}

\[
  x_{n+1}=2c x_n-x_{n-1}
\]
を考える.
これは特性方程式
\[
  r^2-2cr+1=0
\]
をもつ漸化式で,\,$|c|<1$ なら振動的な解が現れる.

ここで
\[
  I_n=x_n^2-2c x_nx_{n-1}+x_{n-1}^2
\]
とおく.
$x_{n+1}=2cx_n-x_{n-1}$ を代入すると,
\[
\begin{aligned}
 I_{n+1}
 &=x_{n+1}^2-2cx_{n+1}x_n+x_n^2\\
 &=(2cx_n-x_{n-1})^2
   -2c(2cx_n-x_{n-1})x_n+x_n^2\\
 &=x_n^2-2cx_nx_{n-1}+x_{n-1}^2\\
 &=I_n.
\end{aligned}
\]
よって $I_n$ は $n$ によらず一定である.

初期値を一つずつ追跡しなくても,
\[
  x_n^2-2c x_nx_{n-1}+x_{n-1}^2=I_0
\]
という曲線上を点 $(x_n,x_{n-1})$ が動くことが分かる.
二次形式の符号や形を調べれば,解が発散するか,振動するかを幾何学的に考えられる.

\subsection{不変量を見つけるための試行}

不変量は,いつも偶然見つかるわけではない.
次のような候補を試す.

\begin{enumerate}
  \item 漸化式が二次式なら,二次式の組合せを候補にする.
  \item $x_{n+1}$ と $x_n$ が対称に現れるなら,対称な式を試す.
  \item 積が現れるなら,積や逆数の和を試す.
  \item 行列で書けるなら,二次形式 $X_n^{\mathsf T}QX_n$ を試す.
\end{enumerate}

行列の形 $X_{n+1}=AX_n$ に対して,
\[
  X_{n+1}^{\mathsf T}QX_{n+1}
  =X_n^{\mathsf T}QX_n
\]
が成り立つためには,
\[
  A^{\mathsf T}QA=Q
\]
であればよい.
これは,回転が長さを保存することの一般化である.
漸化式の不変量探しが,線形代数の問題へ変わった.

\section{発展パターン4：成長のスケールを先に決める}

\subsection{一般項が見えなくても,桁は予想できる}

漸化式の一般項を求めることが難しくても,
$a_n$ が $n$ に対してどの程度の大きさになるかは調べられる.
この見方を漸近解析という.

例えば,
\[
  a_{n+1}=a_n+\frac1{a_n},\qquad a_0>0
\]
を考える.
増分は小さくなるので,\,$a_n$ は直線的には増えなさそうである.
しかし二乗を取ると,
\[
\begin{aligned}
  a_{n+1}^2
  &=\left(a_n+\frac1{a_n}\right)^2\\
  &=a_n^2+2+\frac1{a_n^2}.
\end{aligned}
\]
したがって,
\[
  a_{n+1}^2-a_n^2=2+\frac1{a_n^2}.
\]
右辺はほぼ $2$ なので,\,$a_n^2$ はほぼ $2n$ と考えられる.
つまり
\[
  a_n\text{ は }\sqrt{2n}\text{ 程度}
\]
という予想が得られる.

ここで重要なのは,最初から答えを当てたことではない.
増分の形を見て,どの量を変換すれば一定の増加に近づくかを考えたことである.

\subsection{主項を仮定して,誤差を調べる}

「$a_n$ は $\sqrt{2n}$ くらい」という予想を,
\[
  a_n=\sqrt{2n}+e_n
\]
と書いてみる.
このとき $e_n$ が $\sqrt{n}$ に比べて小さくなるかを調べる.
主項と誤差を分けると,漸化式の本質的な大きさと,
そこからのずれを別々に追跡できる.

一般に,
\[
  a_{n+1}=a_n+g(a_n)
\]
で $g(x)$ が小さくなるとき,\,$a_n$ のスケールを決めるために
\[
  \frac{da}{dn}=g(a)
\]
という連続的な近似を考えることがある.
これは厳密な置き換えではなく,増え方の規模を予想するための道具である.

\subsection{記号 $\sim$ の意味}

二つの正の数列 $a_n,b_n$ に対して,
\[
  a_n\sim b_n
\]
とは,
\[
  \lim_{n\to\infty}\frac{a_n}{b_n}=1
\]
を意味する.
これは「差が小さい」という意味ではない.
例えば $a_n=n+1$ と $b_n=n$ では差は $1$ のままだが,
\[
  \frac{n+1}{n}\longrightarrow1
\]
なので $a_n\sim b_n$ である.
大きさの比が最終的に一致する,という主張である.

\section{発展パターン5：周期的な規則を一周期にまとめる}

\subsection{係数が毎回違うとき,周期を探す}

\[
  a_{n+1}=c_n a_n
\]
で,係数が周期 $p$ をもつとする.
\[
  c_{n+p}=c_n.
\]
一項ずつ見ると係数が変わるが,\,$p$ 回まとめると,
\[
  a_{n+p}
  =c_{n+p-1}\cdots c_{n+1}c_n\,a_n
\]
となる.
周期性により,右辺の積は $n$ の位置によって変わる部分を整理できる.
特に $n=0$ から一周期分をまとめれば,
\[
  a_{mp}=R^m a_0,\qquad
  R=c_{p-1}\cdots c_1c_0
\]
と書ける.

毎回の動きが分からなくても,一周期あたりの倍率 $R$ が分かれば,
長期的な増減を判断できる.

\subsection{周期係数の二階漸化式}

二階の場合も,状態ベクトルを使えば
\[
  X_{n+1}=A_nX_n
\]
である.
係数が周期 $p$ なら,
\[
  X_{(m+1)p}=BX_{mp}
\]
となる.ただし
\[
  B=A_{p-1}\cdots A_1A_0
\]
である.
この $B$ を一周期の遷移行列と見る.

ここでは,一回の規則を解く代わりに,
規則のまとまりを一つの規則として扱った.
「細かい時間では複雑だが,粗い時間では単純」という考え方は,
周期運動や微分方程式の数値計算にも現れる.

\section{発展パターン6：漸化式をアルゴリズムとして読む}

\subsection{式は計算手順でもある}

漸化式
\[
  a_{n+1}=f(a_n)
\]
は,数列の定義であると同時に,計算アルゴリズムでもある.
初期値から出発し,同じ処理を繰り返しているからである.

この見方をすると,次の問いが自然になる.

\begin{itemize}
  \item 計算誤差は何倍に増幅されるか.
  \item 桁数が増えすぎないか.
  \item 同じ計算を何度もせずに済む方法はないか.
  \item $n$ 回の更新を高速にまとめる方法はないか.
\end{itemize}

\subsection{行列の累乗を二乗法で高速化する}

行列漸化式
\[
  X_{n+1}=AX_n
\]
では,素朴には $A$ を $n$ 回掛ける.
しかし,
\[
  A^8=(A^4)^2,\qquad
  A^4=(A^2)^2
\]
のように,指数を二進法で分解すれば,
必要な行列積の回数を大幅に減らせる.

例えば $n=13=8+4+1$ なら,
\[
  A^{13}=A^8A^4A
\]
であり,\,$A,A^2,A^4,A^8$ を順に作ればよい.
これは漸化式の理論が,アルゴリズムの設計へ自然に進む例である.

\subsection{安定性という計算の視点}

真の値を $a_n$,計算で得た値を $\widetilde a_n$ とする.
更新関数が微分可能なら,誤差
\[
  e_n=\widetilde a_n-a_n
\]
は局所的に
\[
  e_{n+1}\simeq f'(a_n)e_n
\]
に従う.
したがって,数学的には収束する漸化式でも,
計算上は誤差が膨らむことがある.

解の公式を持つことと,安定に計算できることは別である.
この違いに気づくと,漸化式は数値解析という分野へつながる.

\section{発展パターン7：刻み幅を細かくして連続世界を見る}

\subsection{一回の変化を小さくする}

漸化式
\[
  x_{n+1}=x_n+h f(x_n)
\]
を考える.ここで $h$ は一回の更新の大きさである.
移項して割ると,
\[
  \frac{x_{n+1}-x_n}{h}=f(x_n)
\]
となる.

もし $h$ を限りなく小さくし,時刻を $t_n=nh$ と考えるなら,
\[
  \frac{x_{n+1}-x_n}{h}
\]
は微分係数に近づく.
その結果,
\[
  \frac{dx}{dt}=f(x)
\]
という微分方程式が現れる.

漸化式が微分方程式に似ているのではない.
一回の変化を小さくする極限を取ろうとしたとき,
差分が微分へ変わるのである.

\subsection{指数関数は,離散と連続の境界にある}

\[
  x_{n+1}=(1+h\lambda)x_n
\]
ならば,
\[
  x_n=(1+h\lambda)^n x_0.
\]
時刻 $t=nh$ を固定して $h\to0$ とすると,\,$n=t/h$ だから,
\[
  (1+h\lambda)^n
  =\left(1+h\lambda\right)^{t/h}
  \longrightarrow e^{\lambda t}.
\]
したがって,指数関数は突然現れるのではなく,
「小さい倍率を非常に多く積み重ねる」という離散的な極限から現れる.

\subsection{安定性の境界が変わる}

連続方程式
\[
  x'=\lambda x,\qquad \lambda<0
\]
では,解は $e^{\lambda t}$ で減衰する.
しかしオイラー型の漸化式は
\[
  x_{n+1}=(1+h\lambda)x_n
\]
であるため,計算が減衰する条件は
\[
  |1+h\lambda|<1
\]
である.
すなわち
\[
  0<h<-\frac2\lambda.
\]

微分方程式の解が安定でも,刻み幅 $h$ を大きくしすぎると,
近似計算は発散する.
連続世界の性質を離散的なアルゴリズムで再現するには,
刻み幅に条件がある.

\section{発展パターン8：解けないことを,失敗ではなく情報にする}

\subsection{一般項を求めるという問いを疑う}

非線形漸化式では,一般項が初等関数で書けないことが多い.
そのとき,線形漸化式の方法を無理に適用しても,
式はますます複雑になるだけである.

そこで問いを変える.

\begin{itemize}
  \item 固定点はどこか.
  \item 増加・減少は保たれるか.
  \item ある区間から外へ出るか.
  \item 周期軌道はあるか.
  \item $n$ が大きいときのスケールは何か.
  \item 数値計算をすれば,どの現象が見えるか.
\end{itemize}

これは「解けないから妥協する」ということではない.
一般項だけを数学の答えだと思わず,規則が生む構造を直接調べるのである.

\subsection{発展の地図}

ここまでのパターンを,問いから整理すると次のようになる.

\begin{center}
\begin{tabular}{c|c}
  問い & 発展する道具\\ \hline
  過去の情報をどう持つか & 状態ベクトル・行列\\
  どこへ近づくか & 固定点・安定性\\
  何が変わらないか & 不変量・二次形式\\
  どのくらい増えるか & 漸近解析・主項\\
  周期的な変化をまとめたい & 一周期の写像・行列積\\
  計算を速く安定にしたい & アルゴリズム・誤差解析\\
  微小な変化の極限を見たい & 微分方程式\\
\end{tabular}
\end{center}

漸化式は,単独の章で閉じる話題ではない.
局所的な更新規則を調べるだけで,
線形代数,解析学,幾何学,数値計算へ道が開く.

\section*{この案を採用する場合の構成上の提案}

この案は,特殊関数を詳しく紹介する前段として置くと流れがよい.
まず「漸化式を見たら,どんな方向へ発展できるか」という方法論を示し,
その後に,実際に何度も現れる特別な解を特殊関数として整理する.

したがって,現行付録に組み込む場合は,
\[
  \text{漸化式の発展パターン}
  \longrightarrow
  \text{特殊関数という到達点}
\]
という二段構成にできる.
前者は思考法,後者はその思考法が生み出した数学的対象を扱う章である.
